\clearpage
\begin{proofbox}
\textbf{方法一｜最低门槛：单调性定义作差。}

第一步：任取同一支上的 $x_1<x_2$，于是 $x_1x_2>0$。

第二步：逐项作差、通分、提取公因式：
\[
\begin{aligned}f(x_2)-f(x_1)
&=Ax_2-B/x_2-Ax_1+B/x_1\\
&=A(x_2-x_1)+B\left(\frac1{x_1}-\frac1{x_2}\right)\\
&=A(x_2-x_1)+\frac{B(x_2-x_1)}{x_1x_2}\\
&=(x_2-x_1)\left(A+\frac B{x_1x_2}\right).
\end{aligned}
\]
第三步：$x_2-x_1>0$，$A>0$，$B/(x_1x_2)>0$，所以差严格大于零。

第四步：由单调性定义，两支分别严格递增。跨过零点时乘积为负，不能沿用第三步。
\end{proofbox}
\clearpage
\begin{proofbox}
\textbf{方法二｜考场最快：同增相加。}

第一步：$Ax$ 在实数范围严格递增，因为 $A>0$。

第二步：反比例函数 $B/x$ 在两个定义区间分别严格递减，因为 $B>0$。

第三步：乘以 $-1$ 后，$-B/x$ 在两个区间分别严格递增。

第四步：两个严格递增函数在同一区间相加仍严格递增，因此得到结论。此法依赖已掌握反比例函数单调性，不需要求导。
\end{proofbox}
\clearpage
\begin{proofbox}
\textbf{方法三｜完整值域：判别式与韦达定理。}

第一步：任取实数 $c$，令 $Ax-B/x=c$，其中 $x\ne0$。

第二步：乘以 $x$ 并移项，得 $Ax^2-cx-B=0$。

第三步：判别式 $\Delta=(-c)^2-4A(-B)=c^2+4AB>0$，因此有两个不同实根。

第四步：由韦达定理，$x_1x_2=-B/A<0$，所以两根都非零且一正一负。原方程的变形可逆，因此两根均满足原式。

第五步：任意实数 $c$ 都可在正、负两支分别且唯一取得，故每支值域都是 $\mathbb R$。求根公式进一步给出
\[
 x_\pm=\frac{c\pm\sqrt{c^2+4AB}}{2A}.
\]
因 $\sqrt{c^2+4AB}>|c|$，故 $x_+>0$，$x_-<0$；两根的和为 $c/A$。
\end{proofbox}
\clearpage
\begin{proofbox}
\textbf{方法四｜常规求导：同时得到变化方向。}

第一步：在 $x\ne0$ 时，$f'(x)=A+B/x^2$。

第二步：由于 $x^2>0$，$A,B>0$，所以 $f'(x)>0$。

第三步：分别在 $(-\infty,0)$ 和 $(0,+\infty)$ 使用导数单调性判定，得到严格递增。不可跨越不属于定义域的 $0$ 使用该判定。

\medskip
\textbf{补充证明：奇偶性与零点。}
定义域关于原点对称，且
\[f(-x)=A(-x)-\frac B{-x}=-Ax+\frac Bx=-f(x).\]
故为奇函数。令 $f(x)=0$，因 $x\ne0$，等价于 $Ax^2=B$，即 $x=\pm\sqrt{B/A}$，两根均符合定义域。
\end{proofbox}
